\documentclass[12pt]{book}
\usepackage{precalc1}
\setcounter{chapter}{7}
\begin{document}

\chapter{Implicit Curves and Conics}

A function assigns one output to each input, and its graph is a curve that passes the
vertical line test. Many important curves in the plane are not of this kind. A circle,
an ellipse, the orbit of a planet, or the set of points where two quantities balance may
be easier to describe by an equation that $x$ and $y$ satisfy together. Such a curve is
\emph{implicit}: the equation states a condition, and the curve is the set of all points
meeting it.

This week broadens the meaning of ``graph.'' Until now, we have mostly drawn graphs by
starting with an input $x$ and computing an output $y$. For implicit curves, we instead
ask which pairs $(x,y)$ satisfy a condition. Sometimes that condition can be solved for
$y$ as a function of $x$. Often it cannot, and that is not a failure. The implicit
equation is then the cleaner description.

\section{Curves Defined by an Equation}

An equation in $x$ and $y$ singles out the points whose coordinates satisfy it. That set
of points is a curve in the plane, whether or not $y$ can be isolated. This changes the
job slightly. Instead of asking ``what is the output when the input is $x$?'' we ask
``does this point belong to the set?''

\begin{definition}{Implicit curve}
An \emph{implicit curve} is the set of all points $(x,y)$ satisfying an equation
$F(x,y)=0$. The equation need not determine $y$ as a function of $x$; a vertical line
may meet the curve in several points.
\end{definition}

The unit circle $x^{2}+y^{2}=1$ is the standard example. No single function describes
it, since most vertical lines through the interior cross it twice. Solving for $y$
yields $y=\pm\sqrt{1-x^{2}}$, two functions, the upper and lower semicircles, that
together reconstruct the curve. The implicit equation holds both halves at once.

\newpage
\begin{keyidea}
An implicit equation describes a curve as a condition its points satisfy, not as a
rule producing $y$ from $x$. When $y$ cannot be isolated as a single function, the
curve simply is not the graph of a function, and the implicit form is the honest
description.
\end{keyidea}

\begin{example}[testing membership]
Does $(3,-4)$ lie on $x^{2}+y^{2}=25$? Substitute: $9+16=25$, true, so the point is
on the curve. Does $(2,3)$? $4+9=13\neq25$, so it is not. Membership is checked by
substitution into the defining equation, never by solving anything.
\end{example}

The example shows the basic habit. An implicit equation is a test. Points that pass the
test are on the curve; points that fail the test are not. A graph is the collection of all
points that pass.

\section{An Elliptic Curve}

A richer example shows how much structure an implicit equation can carry. The point here
is not to develop elliptic curves fully, but to see how a short equation can define a
shape that is not the graph of a function and is not a conic. The equation
\[
y^{2}=x^{3}-x
\]
defines an \emph{elliptic curve}. Reading it as a condition, a point $(x,y)$ lies on
the curve when its $y$-coordinate squared equals $x^{3}-x$. Since a square is never
negative, real points occur only where $x^{3}-x\ge0$, that is where $x(x-1)(x+1)\ge0$:
on the interval $-1\le x\le0$ and on $x\ge1$.

\begin{visual}
The real locus of $y^{2}=x^{3}-x$ falls into two separate pieces: a bounded oval over
$-1\le x\le0$, and an unbounded branch opening to the right over $x\ge1$. Between $0$
and $1$ the right-hand side is negative, so no real point exists, and the curve has a
genuine gap. Such a disconnected shape can never be the graph of a function.
\end{visual}

\begin{center}
\begin{tikzpicture}[xscale=2.0,yscale=1.0]
  \draw[->] (-1.6,0)--(2.2,0) node[right]{\scriptsize $x$};
  \draw[->] (0,-2.4)--(0,2.4) node[above]{\scriptsize $y$};
  \foreach \x in {-1,1,2} \draw (\x,0.05)--(\x,-0.05) node[below]{\scriptsize \x};
  % oval over [-1,0], parametric upper then lower
  \draw[pcblue,thick,smooth,samples=80,domain=-1:0,variable=\x]
    plot (\x,{sqrt(max(0,\x*\x*\x-\x))});
  \draw[pcblue,thick,smooth,samples=80,domain=-1:0,variable=\x]
    plot (\x,{-sqrt(max(0,\x*\x*\x-\x))});
  % right branch over [1,1.85]
  \draw[pcblue,thick,smooth,samples=80,domain=1:1.85,variable=\x]
    plot (\x,{sqrt(max(0,\x*\x*\x-\x))});
  \draw[pcblue,thick,smooth,samples=80,domain=1:1.85,variable=\x]
    plot (\x,{-sqrt(max(0,\x*\x*\x-\x))});
  \node[pcblue] at (-0.5,1.4) {\scriptsize oval};
  \node[pcblue] at (1.7,2.0) {\scriptsize branch};
\end{tikzpicture}
\end{center}

\begin{example}[points on the elliptic curve]
The points with $x=-1$, $x=0$, and $x=1$ all give $y^{2}=0$, so $(-1,0)$, $(0,0)$, and
$(1,0)$ lie on the curve; these are where it meets the $x$-axis. At $x=2$,
$y^{2}=8-2=6$, so $(2,\sqrt6)$ and $(2,-\sqrt6)$ are both on the curve, a vertical
line meeting it twice. At $x=\tfrac12$, $x^{3}-x=\tfrac18-\tfrac12=-\tfrac38<0$, so no
real point sits above $x=\tfrac12$, confirming the gap.
\end{example}

This calculation is the implicit viewpoint in action. We do not need a single formula
$y=f(x)$ to understand the curve. We can read the allowed $x$-values from the condition
$x^{3}-x\ge0$, and then each allowed $x$ gives the possible $y$-values by taking square
roots.

\begin{pitfall}
The symmetry $y\mapsto-y$ in $y^{2}=x^{3}-x$ means every point has a mirror image
across the $x$-axis. This is exactly why the curve is not a function: each admissible
$x$ with $x^{3}-x>0$ produces two $y$-values, one positive and one negative.
\end{pitfall}

\section{Conic Sections}

The most important implicit curves in precalculus are the \emph{conics}, the curves cut
from a cone by a plane and, equivalently, the curves defined by second-degree equations
in $x$ and $y$. Four shapes arise: the circle, the ellipse, the parabola, and the
hyperbola. Each has a standard implicit equation in which the parameters read off
directly as geometric measurements.

Our goal here is recognition and interpretation. We want to look at an equation, name
the kind of curve it represents, and read the main geometric features. This is the same
habit used for lines, parabolas, polynomials, and rational functions: the algebra should
point to the shape.

\begin{keyidea}
\textbf{The conic sections} are the curves of second degree. In standard position,
centered or with vertex at a convenient point:
\begin{itemize}\setlength{\itemsep}{3pt}
  \item \emph{Circle:} $x^{2}+y^{2}=r^{2}$, radius $r$.
  \item \emph{Ellipse:} $\dfrac{x^{2}}{a^{2}}+\dfrac{y^{2}}{b^{2}}=1$, semi-axes $a,b$.
  \item \emph{Parabola:} $y=\dfrac{1}{4p}x^{2}$, or $x^2=4py$, focus distance $p$.
  \item \emph{Hyperbola:} $\dfrac{x^{2}}{a^{2}}-\dfrac{y^{2}}{b^{2}}=1$, opening left
        and right.
\end{itemize}
\end{keyidea}

The parabola has already appeared as a graph of a quadratic function. In conic language,
it is one member of a larger family of second-degree curves. The circle, ellipse, and
hyperbola are usually most naturally described implicitly, because they generally fail
the vertical line test.

\subsection{Circle and ellipse}

The circle is the locus of points at fixed distance $r$ from a center; placing the
center at the origin gives $x^{2}+y^{2}=r^{2}$ from the distance formula. Stretching
the circle by different factors along the two axes produces the ellipse: dividing $x$
by $a$ and $y$ by $b$ before squaring turns the radius-one circle into
$\dfrac{x^{2}}{a^{2}}+\dfrac{y^{2}}{b^{2}}=1$, whose graph reaches $\pm a$ along the
$x$-axis and $\pm b$ along the $y$-axis.

\begin{center}
\begin{tikzpicture}[scale=0.9]
  % circle
  \begin{scope}
    \draw[->] (-1.8,0)--(1.8,0) node[right]{\scriptsize $x$};
    \draw[->] (0,-1.8)--(0,1.8) node[above]{\scriptsize $y$};
    \draw[pcblue,thick] (0,0) circle (1.3);
    \draw[pcred,dashed] (0,0)--(1.3,0) node[midway,below]{\scriptsize $r$};
    \node at (0,-2.3) {\scriptsize $x^2+y^2=r^2$};
  \end{scope}
  % ellipse
  \begin{scope}[xshift=5.5cm]
    \draw[->] (-2.2,0)--(2.2,0) node[right]{\scriptsize $x$};
    \draw[->] (0,-1.8)--(0,1.8) node[above]{\scriptsize $y$};
    \draw[pcblue,thick] (0,0) ellipse (1.8 and 1.1);
    \draw[pcred,dashed] (0,0)--(1.8,0) node[midway,below]{\scriptsize $a$};
    \draw[pcred,dashed] (0,0)--(0,1.1) node[midway,left]{\scriptsize $b$};
    \node at (0,-2.3) {\scriptsize $\frac{x^2}{a^2}+\frac{y^2}{b^2}=1$};
  \end{scope}
\end{tikzpicture}
\end{center}

The denominators in the ellipse equation measure squared distances to the edge along
the coordinate axes. A larger denominator lets the curve stretch farther in that
direction. This is why the square roots of the denominators, not the denominators
themselves, are the semi-axis lengths.

\begin{example}[reading an ellipse]
The equation $\dfrac{x^{2}}{9}+\dfrac{y^{2}}{4}=1$ has $a^{2}=9$ and $b^{2}=4$, so
$a=3$ and $b=2$. The ellipse stretches to $\pm3$ along the $x$-axis and $\pm2$ along
the $y$-axis. Setting $y=0$ gives $x=\pm3$; setting $x=0$ gives $y=\pm2$, confirming
the intercepts.
\end{example}

\subsection{Hyperbola}

Changing the plus sign in the ellipse to a minus produces the hyperbola
$\dfrac{x^{2}}{a^{2}}-\dfrac{y^{2}}{b^{2}}=1$. The curve has two branches opening left
and right. Unlike an ellipse, it does not close up; its branches continue forever while
approaching diagonal lines called asymptotes.

For large $x$, the constant $1$ becomes negligible and the equation approaches
$\dfrac{x^{2}}{a^{2}}=\dfrac{y^{2}}{b^{2}}$. Solving that approximate equation gives the
asymptotes $y=\pm\dfrac{b}{a}x$.

\begin{center}
\begin{tikzpicture}[scale=0.8]
  \draw[->] (-3.4,0)--(3.4,0) node[right]{\scriptsize $x$};
  \draw[->] (0,-2.8)--(0,2.8) node[above]{\scriptsize $y$};
  % asymptotes y = +- (b/a) x, take a=1.4,b=1.6 -> slope ~1.14
  \draw[pcgreen,dashed] (-2.6,-2.97)--(2.6,2.97);
  \draw[pcgreen,dashed] (-2.6,2.97)--(2.6,-2.97);
  \node[pcgreen] at (2.6,2.3) {\scriptsize $y=\frac{b}{a}x$};
  % right branch x = a cosh t, y = b sinh t
  \draw[pcblue,thick,smooth,samples=60,domain=-1.5:1.5,variable=\t]
    plot ({1.4*cosh(\t)},{1.6*sinh(\t)});
  \draw[pcblue,thick,smooth,samples=60,domain=-1.5:1.5,variable=\t]
    plot ({-1.4*cosh(\t)},{1.6*sinh(\t)});
  \fill (1.4,0) circle (1.6pt) node[below right]{\scriptsize $a$};
\end{tikzpicture}
\end{center}

\begin{example}[asymptotes of a hyperbola]
For $\dfrac{x^{2}}{16}-\dfrac{y^{2}}{9}=1$, $a^{2}=16$ and $b^{2}=9$, so $a=4$, $b=3$.
The vertices sit at $(\pm4,0)$ and the asymptotes are $y=\pm\tfrac34 x$. The branches
hug those lines as $x$ grows.
\end{example}

The vertices come from setting $y=0$, just as the ellipse intercepts did. The
asymptotes come from the end behavior. A good sketch of a hyperbola should show both:
where the branches begin and which lines they approach.

\subsection{Recognizing a conic from its equation}

A second-degree equation with no $xy$ term is a conic, and the signs of the squared
terms name which one. Equal positive coefficients on $x^{2}$ and $y^{2}$ give a
circle; unequal positive coefficients give an ellipse; opposite signs give a
hyperbola; a single squared term gives a parabola.

\begin{keyidea}
\textbf{Identifying a conic} from $Ax^{2}+Cy^{2}+\cdots=0$ (no $xy$ term):
\begin{itemize}\setlength{\itemsep}{2pt}
  \item $A=C$ (same sign): circle.
  \item $A\neq C$, same sign: ellipse.
  \item $A,C$ opposite signs: hyperbola.
  \item exactly one of $A,C$ is zero: parabola.
\end{itemize}
\end{keyidea}

\begin{example}[classifying by inspection]
$4x^{2}+4y^{2}=36$ has equal coefficients, a circle (radius $3$ after dividing).
$x^{2}+9y^{2}=9$ has unequal positive coefficients, an ellipse. $x^{2}-y^{2}=1$ has
opposite signs, a hyperbola. $y-x^{2}=0$ has a single squared term, a parabola.
\end{example}

After identifying the type, the next step is usually to rewrite the equation into a
standard form. Standard form exposes the radius, axes, vertices, or asymptotes. The
classification tells us what kind of curve to expect; the standard form tells us its
measurements.

\section{Exercises}

\begin{exercises}
\item Decide by substitution whether each point lies on $x^{2}+y^{2}=25$:
(a) $(0,5)$; (b) $(3,4)$; (c) $(-4,3)$; (d) $(2,2)$.

\item Solve $x^{2}+y^{2}=9$ for $y$, obtaining two functions, and state the domain of
each. Explain how they recombine to the full circle.

\item For the elliptic curve $y^{2}=x^{3}-x$, find all points with $x=-1,0,1$, and
determine whether any real point has $x=\tfrac14$.

\item Explain, using the symmetry $y\mapsto-y$, why $y^{2}=x^{3}-x$ is not the graph
of a function, and describe its two pieces.

\item Find every point on $y^{2}=x^{3}-x$ with $x=3$, and verify each by substitution.

\item Write the equation of the circle centered at the origin passing through
$(6,8)$, and state its radius.

\item Identify $a$ and $b$ for $\dfrac{x^{2}}{25}+\dfrac{y^{2}}{16}=1$, give the four
axis intercepts, and sketch the ellipse.

\item For $\dfrac{x^{2}}{9}-\dfrac{y^{2}}{4}=1$, give the vertices and the equations of
the asymptotes, and sketch both branches with the asymptotes.

\item Classify each conic from its equation: (a) $3x^{2}+3y^{2}=12$;
(b) $4x^{2}+y^{2}=16$; (c) $y^{2}-x^{2}=1$; (d) $x-2y^{2}=0$.

\item Explain why the hyperbola $\dfrac{x^{2}}{a^{2}}-\dfrac{y^{2}}{b^{2}}=1$
approaches the lines $y=\pm\dfrac{b}{a}x$ for large $x$, reasoning from the equation.

\item Starting from the unit circle $x^{2}+y^{2}=1$, explain how dividing $x$ by $a$
and $y$ by $b$ before squaring produces an ellipse, and what $a$ and $b$ measure on
the resulting curve.

\item A second-degree equation $Ax^{2}+Cy^{2}=1$ has $A=4$, $C=-1$. Name the conic,
give its vertices, and state its asymptotes.
\end{exercises}

\end{document}
